% ECONOMICS_PROBLEM: {"id": "D81-16", "title": "Separate identification of beliefs and preferences", "jel_code": "D81", "source_id": "OP-0587"}
% FIRST_POSTED: 2026-10-11
% ATTEMPT_STATUS: solved
% ENTRY_KIND: research_attempt
% SUBMISSION: {"kind": "research_attempt", "category": "economics", "problem_id": "D81-16", "model": {"provider": "Anthropic", "version": "Claude Fable 5.1 (claude-fable-5-1)", "reasoning": "Claude Code agent session, default effort", "generated_on": "2026-10-11"}, "other_models": "None for the Lean development; the hand proof that is formalized was generated by OpenAI Codex (pull request #162, manuscript linked in the text)", "tools": "Lean 4.34.1 with Mathlib d13f23b723b8a846827a245b89c10fc7d3f11612 and lake; compilation on a separate Linux server", "target": "The D81-16 target as resolved by the hand proof for its declared model classes: the finite-consequence logit choice model with invariant normalized utilities and beliefs (informative designs, sharp equivalence classes, identified policy functionals, minimal information dimension) and the calibrated max-min extension with compact convex prior sets (compatible prior family, sharp and joint policy values, whole-prior-set identification)", "scope": "Theorems 1, 2, 3, 4 and Proposition 1 of the hand proof are formalized in full for all finite m, r >= 1, boundary priors, incompatible observations, J = 0 and L = 0, known and unknown scale; the polytope property of the witness projection is formalized by Fourier-Motzkin elimination; not formalized are the two algorithmic citations (a real quantifier-elimination procedure deciding the recognition system, and the bit complexity of rational linear programming), which concern the running of algorithms, and the worked examples", "human_contributions": "None; the submitter specified the task and submitted the result, with no mathematical edits", "verification": "Lean kernel check of every statement with only the standard axioms (propext, Classical.choice, Quot.sound); model self-audit of the definitions against the manuscript; no independent human proof verification", "reproducibility": "Public Lean repository pinned to a commit, with build commands and the actual checking outcome reported; the interactive agent-session prompts are not published"}
% FULL_SOLUTION_SCOPE: {"target": "The D81-16 target as formulated by the hand proof for its declared model classes (Theorems 1-4 and Proposition 1 of the manuscript)", "result": "Ident/Main.lean: theorem1_bijection, theorem1_fiber, theorem1_secant (normalized matrix bijection, sharp fibers, secant criterion and recognition system); theorem2 (the least number of independent log odds of a globally identifying experiment is m + r - 2, and m + r - 1 with unknown scale); theorem3, theorem3_model (compatibility criterion and the sharp convex-prior family, in data and observation-law terms); theorem4_scalar, theorem4_joint (sharp policy-value interval, joint identified set as the projection of the finite linear witness system, compact and convex, hull witnesses, attained contrast extrema); proposition1 (vertex criterion, identified set equals H); Ident/Polyhedron.lean and Ident/Polytope.lean: Fourier-Motzkin elimination (linear images of polyhedra are polyhedra) and the polytope statements isPolytope_Hset, isPolytope_Fset, isPolytope_jointSet", "coverage": "All finite m, r >= 1 (m = r = 1 included), boundary priors, every finite experiment (finitely many finite menus with reference actions and consequence labels), arbitrary positive known scale and unknown scale, incompatible observed laws (empty fibers), J = 0 and L = 0 in the ambiguity extension, lower-dimensional H; both directions of every equivalence", "dependencies": "Mathlib only (d13f23b7): the inverse function theorem (C^1 local inverse), rank-nullity, Krein-Milman, compactness of finite products, intermediate value theorem on convex sets; no additional axioms", "human_verification": "None", "unverified": "The transcription of the manuscript's definitions into Lean was checked by the model only; the hand proof itself is an LLM claim pending community expert review; the existence of a terminating decision procedure for the recognition system (real quantifier elimination) and the polynomial bit complexity of rational linear programming are algorithmic citations used by the manuscript to describe computation and are not formalized (the mathematical statements those procedures decide are); the worked examples are not formalized", "novelty": "Not checked beyond the manuscript's own references (Manski 2018; Basu, Pollack and Roy 2006; Rudin 1976; Ben-Tal and Nemirovski 2023)"}
\documentclass[11pt]{article}
\usepackage[margin=1in]{geometry}
\usepackage{amsmath,amssymb,amsthm}
\usepackage[hidelinks]{hyperref}
\newtheorem{theorem}{Theorem}
\newtheorem{proposition}{Proposition}
\newcommand{\R}{\mathbb R}
\newcommand{\M}{\mathcal M}
\newcommand{\ind}{\mathbf 1}
\DeclareMathOperator{\conv}{conv}
\DeclareMathOperator{\vecop}{vec}
\title{A complete machine-checked proof of the finite-experiment separation of beliefs, preferences and ambiguity}
\author{Claude Fable 5.1 (Claude Code), submitted by Ji Ho Bae}
\date{October 11, 2026}
\begin{document}
\maketitle
\begin{abstract}
The hand proof for this problem resolves the informative-model-class target of D81-16 for a declared finite-consequence logit choice model and its calibrated max--min extension: a structural criterion and sharp fibers for the separate identification of beliefs and normalized utility, the minimal information dimension $m+r-2$ (one more with unknown scale), the entire compatible convex-prior family, sharp and joint policy values, and a vertex criterion for identification of the whole prior set. This note reports a complete Lean~4 formalization of that proof against Mathlib: Theorems~1--4 and Proposition~1 are proved for all finite $m,r\ge1$, including boundary priors, incompatible observations, the unknown-scale case and $J=0$. The lower bound of Theorem~2 is derived from Mathlib's inverse function theorem. The development has no \texttt{sorry}, no warnings, assumes nothing beyond Mathlib, and uses only the standard axioms.
\end{abstract}

\section{The statements}

There are $m\ge1$ states, consequences $c_0<\dots<c_r$ with $r\ge1$, utilities $u_0=0<u_1<\dots<u_r=1$ and a prior $p\in\Delta_m$ (boundary priors included). A finite experiment supplies finitely many menus $A_z$ with reference actions $a^0_z$ and known consequence labels $j_z(a,s)$; the observation is the population logit law $\pi_z(a)=\exp(\lambda V_{p,u}(a,z))/\sum_b\exp(\lambda V_{p,u}(b,z))$ with $V_{p,u}(a,z)=\sum_sp_su_{j_z(a,s)}$ and known $\lambda>0$, and $\ell_{z,a}=\lambda^{-1}\log(\pi_z(a)/\pi_z(a^0_z))$ are the $N=\sum_z(|A_z|-1)$ independent log odds. $\M$ is the set of matrices $W\ge0$ with last column summing to one, strictly increasing column sums and $W_{sj}=W_{sr}\sum_tW_{tj}$, and $D$ the integer matrix $D_{(z,a),(s,j)}=\ind\{j_z(a,s)=j\}-\ind\{j_z(a^0_z,s)=j\}$.

\begin{theorem}[Theorem 1]
$W_{sj}=p_su_j$ is a bijection from the normalized family onto $\M$, with inverse $p_s=W_{sr}$, $u_j=\sum_sW_{sj}$. The sharp equivalence class at an observed law is $\mathcal F(\ell)=\{W\in\M:D\vecop W=\ell\}$, an empty fiber reporting incompatibility. The experiment is globally identifying iff $\ker D\cap(\M-\M)=\{0\}$, iff the system $W,\widetilde W\in\M$, $D\vecop(W-\widetilde W)=0$, $\sum(W_{sj}-\widetilde W_{sj})^2>0$ is infeasible.
\end{theorem}

\begin{theorem}[Theorem 2]
The smallest number of independent log odds of a globally identifying finite experiment is $d=m+r-2$; with unknown positive scale, joint identification of $(\lambda,p,u)$ requires and admits $d+1$.
\end{theorem}

In the calibrated max--min model the prior is replaced by a nonempty compact convex $K\subseteq\Delta_m$ and values are $V_K(a)=\min_{p\in K}\sum_sp_su_{j(a,s)}$; sure acts identify every $u_j$ and the endpoint comparison identifies $\lambda$. With $d_j$ the state-utility vectors and $v_j$ the observed values of the acts $a_j$ ($j\in J$), $H=\{p\in\Delta_m:p\cdot d_j\ge v_j\ \forall j\}$ and $F_j=H\cap\{p\cdot d_j=v_j\}$.

\begin{theorem}[Theorem 3]
The data are compatible iff $H$ and every $F_j$ are nonempty, and the sharp family is $\mathcal K(v)=\{K\subseteq H\text{ nonempty compact convex}:K\cap F_j\neq\varnothing\ \forall j\}$.
\end{theorem}

\begin{theorem}[Theorem 4]
For compatible data with $J\ge1$ the sharp set of $\min_{p\in K}p\cdot h$ is $[\min_Hp\cdot h,\ \min_j\max_{F_j}p\cdot h]$. For policies $h_1,\dots,h_L$ the joint identified set of lower utilities is the projection of the finite linear witness system (face witnesses $p^j\in H$, policy witnesses $q^\ell\in H$, $p^j\cdot d_j=v_j$, $q^\ell\cdot h_\ell=w_\ell$, $p^j\cdot h_k\ge w_k$, $q^\ell\cdot h_k\ge w_k$); it is compact and convex, every retained vector has a witnessing hull of at most $J+L$ points, and every linear contrast has attained sharp extrema.
\end{theorem}

\begin{proposition}[Proposition 1]
For compatible data with $J\ge1$, the whole convex prior set is identified iff every vertex of $H$ is the unique point of some $F_j$; it then equals $H$.
\end{proposition}

\section{What is verified}

The Lean development is public at a fixed commit \cite{Lean}. The normalized family is \texttt{Param m r}, an experiment is \texttt{Experiment m r} (environments, finite menus, reference actions, labels), the law is \texttt{Experiment.obs}, the log odds \texttt{Experiment.logOdds}, the matrix family \texttt{Mset}, the matrix $D$ \texttt{Experiment.Dvec}/\texttt{Dlin}, and global identification \texttt{Experiment.Identifies} (injectivity of the law on the family). Prior sets are \texttt{KAdm}, lower values \texttt{VK}, the polytope and faces \texttt{Hset}, \texttt{Fset}, compatibility \texttt{Compatible} and the family \texttt{SharpFamily}; the max--min logit law is \texttt{AmbExperiment.obs}.

\subsection*{The collected statements (\texttt{Ident/Main.lean})}
\begin{verbatim}
theorem theorem1_secant (E : Experiment m r) (hm : 0 < m) (hr : 0 < r) (hlam : 0 < lam) :
    (E.Identifies lam <-> ker E.Dlin  cap (Mset m r - Mset m r) = {0}) /\
    (E.Identifies lam <-> not exists W W', W in Mset m r /\ W' in Mset m r /\
      E.Dvec (W - W') = 0 /\ 0 < sum s, sum j, (W s j - W' s j) ^ 2)

theorem theorem2 (m' r' : Nat) (hlam : 0 < lam) :
    IsLeast {N | exists E : Experiment (m'+1) (r'+1), E.Identifies lam /\ E.obsCard = N}
      (m' + r') /\
    IsLeast {N | exists E : Experiment (m'+1) (r'+1), E.IdentifiesJoint /\ E.obsCard = N}
      (m' + r' + 1)

theorem theorem3 (d : I -> Fin m -> R) (v : I -> R) :
    ((exists K, Compatible d v K) <-> (Hset d v).Nonempty /\ forall j, (Fset d v j).Nonempty) /\
    forall K, Compatible d v K <-> K in SharpFamily d v

theorem theorem4_scalar [Nonempty I] (d v) (hH : (Hset d v).Nonempty)
    (hF : forall j, (Fset d v j).Nonempty) (h : Fin m -> R) :
    exists alpha beta, IsLeast (ip . h '' Hset d v) alpha /\
      (forall j, IsGreatest (ip . h '' Fset d v j) (beta j)) /\
      {w | exists K, Compatible d v K /\ VK K h = w} = {w | alpha <= w /\ forall j, w <= beta j}

theorem theorem4_joint (hm : 0 < m) (d v) (h : L -> Fin m -> R) :
    JointSet d v h = (fun x => x.2.2) '' WitnessSys d v h /\
      IsCompact (JointSet d v h) /\ Convex R (JointSet d v h) /\ (hull witnesses) /\
      forall c, (JointSet d v h).Nonempty ->
        exists a b, (fun w => sum l, c l * w l) '' JointSet d v h = Icc a b

theorem proposition1 [Nonempty I] (d v) (hF : forall j, (Fset d v j).Nonempty) :
    ((forall K K', Compatible d v K -> Compatible d v K' -> K = K') <->
      forall x in (Hset d v).extremePoints R, exists j, Fset d v j = {x}) /\
    ((forall K K', Compatible d v K -> Compatible d v K' -> K = K') ->
      forall K, Compatible d v K -> K = Hset d v)
\end{verbatim}
Unicode is transliterated and some binders abbreviated; \texttt{theorem1\_bijection} and \texttt{theorem1\_fiber} state the bijection with its inverse and the sharp fiber at an arbitrary positive observed law. The manuscript's $m,r\ge1$ are written $m'+1$, $r'+1$ in Theorem~2, so $d=m'+r'$. \texttt{theorem3\_model} is Theorem~3 in observation-law terms: the admissible prior sets observationally equivalent to $K$ in the calibrated max--min model are exactly $\mathcal K(v)$ for the transformed data, after \texttt{AmbExperiment.obs\_eq\_iff} (sure acts identify $u$) and \texttt{obs\_eq\_iff\_raw} (the endpoint comparison identifies $\lambda$).

\section{Architecture of the formal proof}

\begin{itemize}
\item \emph{Logit law} (\texttt{Ident/Logit.lean}). Two value vectors give the same softmax law iff their differences to the reference agree (\texttt{softmax\_eq\_iff}), and a positive normalized vector is determined by its ratios (\texttt{eq\_of\_ratio\_eq}); hence observational equivalence is equality of log odds and the sharp class at an arbitrary law is the fiber of its transform.
\item \emph{Theorem 1} (\texttt{Ident/Model.lean}, \texttt{Ident/Theorem1.lean}). The bijection onto $\M$ is explicit (\texttt{toMatrix\_bijOn}, \texttt{ofMatrix}), zero prior coordinates included; the choice formula gives $D\vecop W=\ell$ (\texttt{logOdds\_eq\_Dvec}); the secant criterion and the recognition system are equivalent reformulations of injectivity of $D$ on $\M$.
\item \emph{Theorem 2} (\texttt{Ident/Dimension.lean}, \texttt{Ident/Rank.lean}). The design with a bet on each state $s<m$ and each sure act $0<j<r$ against $c_0$ recovers $p_s$ and $u_j$ from its $d$ log odds, at boundary priors too; the endpoint comparison adds $\lambda$. For the lower bound, the manuscript's elementary fact (a $C^1$ map from a nonempty open subset of $\R^d$ into $\R^N$, $N<d$, is not injective) is proved as \texttt{not\_injOn\_of\_card\_lt}: at a point of maximal derivative rank $\rho$ the map is straightened by Mathlib's \texttt{ContDiffAt.localInverse}, the rank bound forces the derivative to vanish on the $d-\rho>0$ kernel directions, and the map is constant along a short segment, giving a collision. It is applied to the polynomial log-odds map on the open set of strictly positive priors and strictly ordered utilities (and to $\lambda\,D\vecop(pu^\top)$ for the unknown scale).
\item \emph{Polytopes} (\texttt{Ident/Polyhedron.lean}, \texttt{Ident/Polytope.lean}). A polyhedron is the solution set of finitely many linear inequalities; Fourier--Motzkin elimination of one variable (\texttt{IsPolyhedron.image\_fst\_prod\_real}) gives that linear images of polyhedra in finite-dimensional spaces are polyhedra (\texttt{IsPolyhedron.image}). Hence $H$, the faces, the witness system and the joint identified set are polyhedra, and the compact ones polytopes (\texttt{isPolytope\_jointSet}), as Theorem~4 states.
\item \emph{Theorems 3--4} (\texttt{Ident/Simplex.lean}, \texttt{Ident/Ambiguity.lean}, \texttt{Ident/Policy.lean}, \texttt{Ident/AmbModel.lean}). Lower values are attained by compactness; every compatible set lies in $H$ and meets every face, and conversely the hull of face witnesses (or $H$ itself) is compatible. The interval of Theorem~4 uses the intermediate value theorem on the convex set $H$ to place a point with the required value into the hull of face maximizers; the joint set is shown equal to the projection of the witness system in both directions, closed and bounded (hence compact) and convex, with a singleton prior as witness when $J=L=0$; contrasts attain their extrema by compactness and fill the interval by connectedness.
\item \emph{Proposition 1} (\texttt{Ident/Vertex.lean}). Forced vertices lie in every compatible set, and Mathlib's Krein--Milman theorem (\texttt{closure\_convexHull\_extremePoints}) gives $H\subseteq K$; an unforced vertex is omitted by the hull of other face points because extreme points of a hull belong to the generating set (\texttt{extremePoints\_convexHull\_subset}).
\end{itemize}

\section{Scope}

The formalization follows the hand proof. The matrix $W=pu^\top$ carries its zero column $j=0$, in bijection with the manuscript's $m\times r$ matrices. ``Independent log odds'' are the coordinates $(z,a)$ with $a\ne a^0_z$. ``Vertex of $H$'' is formalized as extreme point, which is the same for the polytope $H$. The joint identified set is proved to be the projection of the finite linear witness system, compact, convex and a polytope (a bounded polyhedron, by the Fourier--Motzkin elimination formalized in \texttt{Ident/Polyhedron.lean}), with hull witnesses of at most $J+L$ points and attained contrast extrema. Two citations of the manuscript are algorithmic and are not formalized: that a real quantifier-elimination procedure \cite{BPR} decides the recognition system in finitely many steps, and the polynomial bit complexity of rational linear programming \cite{BTN}. Both concern the running of an algorithm rather than a mathematical statement about the model; the statements those procedures decide are formalized (\texttt{identifies\_iff\_no\_collision}, \texttt{contrast\_image\_eq\_Icc}, \texttt{isPolytope\_jointSet}). The worked examples of the manuscript are not formalized. Theorem~2's lower bound uses Mathlib's inverse function theorem in place of \cite{Rudin}.

\section{Build and audit}

Lean \texttt{v4.34.1}, Mathlib commit \texttt{d13f23b723b8a846827a245b89c10fc7d3f11612}; about 2,750 lines in 13 files. The commands \texttt{lake exe cache get}, \texttt{lake build} and \texttt{lake env lean Verify.lean} were run on a fresh clone on a separate Linux machine. The build reports no errors and no warnings; the sources contain no \texttt{sorry}. Each of the 54 statements audited in \texttt{Verify.lean}, including the nine collected statements of \texttt{Ident/Main.lean}, depends only on \texttt{propext}, \texttt{Classical.choice} and \texttt{Quot.sound}.

\section{Completion Estimate}
\noindent\textbf{COMPLETION: 100\%}

This is the claim for the catalogue target as resolved by the hand proof for its declared model classes: informative designs and sharp equivalence classes for the finite logit model, the minimal information dimension, and the calibrated max--min extension with its compatible prior family, policy values and whole-set identification. The machine check covers Theorems~1--4 and Proposition~1 for all finite $m,r\ge1$ and all the boundary cases listed above, with no external theorem beyond Mathlib. It is not an independent expert review, and the hand proof itself remains an LLM claim pending community review.

\section*{Attribution and reproduction}
Model and process: Claude Fable 5.1 (Claude Code agent session, Anthropic), 11 October 2026, working from the manuscript \cite{Hand} and the Mathlib library; no other model contributed to the Lean development. Human contributions: none; the submitter specified the task and submitted the result, with no mathematical edits. The complete generated output is the Lean repository at its pinned commit \cite{Lean}; the interactive session prompts are not published. Lean source: \url{https://github.com/jbaelaw/belief-preference-identification-lean-20261011/tree/a4b4cbff0dd12ae9d7ac22f1b833b8e25df9153f} (\texttt{Ident/Main.lean}). Reproduction: Lean \texttt{v4.34.1}, Mathlib commit \texttt{d13f23b723b8a846827a245b89c10fc7d3f11612}, built with \texttt{lake} as described above; the project has several files and depends on Mathlib, so it is not a single-file Lean Web example. Checking outcome: no errors, no warnings, no \texttt{sorry}, standard axioms only, on a fresh clone. The repository maintainers have not verified the result.

\section*{Final status}
SOLVED, as a machine-checked formalization of the hand proof's complete resolution of the declared model-class target; community expert review of the hand proof is pending.

\begin{thebibliography}{9}
\bibitem{Hand} OpenAI Codex (GPT-6 family), submitted by Ji Ho Bae, \emph{Finite experiments separating beliefs, preferences and ambiguity}, hand proof for D81-16, \url{https://github.com/jbaelaw/belief-preference-identification-proof/blob/ce95354f5d52d298d0d3640f9af1de86bc36b266/paper/Proof.tex}.
\bibitem{Lean} Lean 4 formalization, \url{https://github.com/jbaelaw/belief-preference-identification-lean-20261011/tree/a4b4cbff0dd12ae9d7ac22f1b833b8e25df9153f}.
\bibitem{BPR} Saugata Basu, Richard Pollack and Marie-Fran\c{c}oise Roy, \emph{Algorithms in Real Algebraic Geometry}, second edition, Springer, 2006.
\bibitem{Rudin} Walter Rudin, \emph{Principles of Mathematical Analysis}, third edition, McGraw--Hill, 1976, Theorem 9.24.
\bibitem{BTN} Aharon Ben-Tal and Arkadi Nemirovski, \emph{Lecture Notes: Optimization III}, ISYE 6663, Spring 2023, Section 8.4.2.
\end{thebibliography}
\end{document}
